\chapter{Mathematical Model}

\section{Robust communication QoS}
For an uncertain channel $\vect{h}_k=\hat{\vect{h}}_k+\Delta\vect{h}_k$ with
$\|\Delta\vect{h}_k\|_2\leq\epsilon_h$, the worst-case communication SINR
requirement is imposed over the entire uncertainty ball. In covariance form,
the desired communication covariance must dominate multiuser interference,
dedicated sensing interference, and noise for every admissible channel error.
This is a semi-infinite constraint before the S-procedure is applied.

\section{Sensing SINR and PCRB}
The dedicated covariance $\mat{S}_p$ is used to guarantee target detection
through a sensing-SINR threshold. Tracking accuracy is controlled by the data
FIM $\mat{J}_p$ and the trace-PCRB condition
\[
 \operatorname{tr}(\mat{J}_p^{-1})\leq\Gamma_{\rm track,p}.
\]
For the adopted complex Gaussian mean model, the data FIM contains the factor
$2/\sigma_s^2$. The physical derivative matrix $\mat{D}_p$ retains the target
geometry. A small numerical regularization may be used only to prevent singular
linear algebra; it must not be misrepresented as new sensing information.

\section{Original problem P1}
The original formulation jointly minimizes total transmit power over
$\{\vect{w}_k\}$, $\{\mat{S}_p\}$, and $\{b_{mp}\}$ subject to robust
communication QoS, sensing SINR, PCRB, per-AP power, sensing authorization, and
cardinality constraints. The coupled sources of nonconvexity are: beamformer
quadratic products, rank-one covariance relations, worst-case CSI constraints,
PCRB inverse, SINR ratios, and binary variables.

\section{Interpretation of PCRB geometry}
Sensing SINR is largely an energy measure, whereas PCRB depends on the
eigenstructure of $\mat{J}_p$. APs with similar viewing geometry can contribute
energy without providing a second informative direction. This explains why a
distance-only association can be poor for tracking, and why FIM-aware topology
construction is a strong baseline.

\sourcebox{Full derivation reference: \texttt{PROBLEM\_FORMULATION.md},
Sections 2--6; numerical implementation: \texttt{solve\_p3\_sca\_t.m}.}
